\documentclass[11pt]{article}
\usepackage[margin=1in]{geometry}
\usepackage{booktabs}
\usepackage{url}
\title{Perception-radius terminal set without fallback algorithms}
\date{Prepared October 4, 2026}
\begin{document}
\maketitle

\section{Changes (user design)}
\paragraph{No fallback.} Each frame builds one anchor: a potential-field rollout
at startup, otherwise the shifted previous plan extended by path guidance. On
that anchor the controller minimizes PCBF slack, then CLF slack. If any step
fails, the controller reports \texttt{noOptimizationSolution}. Four mechanisms
were removed (the fresh-seed re-solve was later restored, see Addendum): the fresh-seed retry, the input-box enlargement, the inherited and
observer-conditioned slack budgets with their restoration solves, and the
replacement of an unusable shift. Prefix collision rows remain soft, at the
user's choice.

\paragraph{Terminal set.} The free-pose cruise core and its module
(\path{terminalContinuation.m}, 582 lines) are deleted. At the last node the
target estimate must be beyond the perception radius $R=50$~m and separating,
and every ego rectangle corner must lie inside the road,
$\texttt{lateralClearance}=[4;4]$~m. The target rows are linearized with the
anchor direction~$n$: $n^\top(p-q)\ge R$ implies $\|p-q\|\ge R$, and
$n^\top(R(\psi)v-w)\ge 0$ is the separating speed. The ego generator support is
subtracted from every terminal row. The horizon is variable: the anchor stops at
the first node, between \texttt{horizonSteps} and 512, that is separating and at
least $R+1$~m from the target. Without a target, the shortest horizon is used.

\paragraph{Corrections found during validation.}
\begin{itemize}
\item \textbf{Target uncertainty in the terminal rows.} With the target tube
  included, the first noisy frame needed a distance of $50+60.6$~m, because the
  tube's position radius is 60.6~m at 5~s. Every noisy case therefore failed at
  frame~0. The target now enters the terminal rows as its estimate; its
  uncertainty is handled by the collision rows inside~$R$.
\item \textbf{Slack-cap tolerance.} The CLF stage stalled in Clarabel
  (InsufficientProgress) with a $10^{-5}$ slack cap. The cap is now
  $10^{-4}$~m, and four failing holds solve.
\item \textbf{Horizon end.} With the anchor exactly at $R$, the terminal row was
  active at zero correction. Ending the horizon 1~m beyond~$R$ solved six
  stalled exact and noisy holds.
\end{itemize}

\section{Results (final source)}
Of 655 tests, 655 pass. The table compares the campaigns with the earlier
records of this day (plan-innovation trust, cruise terminal core, with fallbacks).

\begin{center}
\begin{tabular}{lccc}
\toprule
 & recovered & no solution & collision\\
\midrule
Exact observations, before & 13 & 1 & 0\\
Exact observations, now & 11 & 3 & 0\\
Noisy estimator, before & 7 & 6 & 1\\
Noisy estimator, now & 1 & 13 & 0\\
\bottomrule
\end{tabular}
\end{center}

With exact observations, the stops are the CLF-stage numerical failures of
8 and 15~m/s brakingLead (holds 26 and 110) and a PCBF infeasibility of 15~m/s
curvedHeadOn at hold 1084, during recovery.

With the noisy estimator, 13 cases stop, mostly in the first 0--65 holds. At
those holds the conic problem is primal infeasible (status~2), not stalled. A
replay of six of them that drops one constraint group at a time gives:
\begin{itemize}
\item \textbf{Road row only.} 15-m/s crossing (2.10~s) and 8-m/s
  acceleratingHeadOn (3.25~s) become solvable when only the terminal road rows
  are removed. The horizons are short (46 and 15 nodes): the target exits~$R$
  while the ego is still offset from the road centre.
\item \textbf{Target or road row.} 15-m/s headOn (1.50~s) becomes solvable when
  either the terminal target rows or the terminal road rows are removed.
\item \textbf{Target row only.} 8-m/s brakingLead (frame~0) becomes solvable
  only without the terminal target rows. Its horizon reached 512 nodes
  (25.6~s) without exiting~$R$: following a lead vehicle never reaches this
  terminal set.
\item \textbf{Collision rows only.} 8-m/s turningCrossing (0.40~s) becomes
  solvable only without the observer-tightened collision rows.
\item \textbf{No single group.} 8-m/s headOn (0.80~s) is not solvable with any
  one group removed.
\end{itemize}
No collision occurred in any executed hold.

\section{Scope}
No invariance or recursive-feasibility claim is made for this terminal set. A
target whose constant sideslip turns it back can re-enter~$R$ after the endpoint.
The terminal set excludes following a lead vehicle inside~$R$. With a short
horizon, the road rows require the ego to be inside the road almost
immediately after the target exits~$R$.

\section{Files}
\path{PERCEPTION_TERMINAL_SET_20261004/} contains per-case metrics, the test
table, the source hashes of the final runs, the constraint-drop replays of the
noisy failures, and the replay script. Raw traces are in
\path{~/.cache/collisionAvoidance/perception-terminal-20261004c/}. Earlier
iterations are in the \path{-20261004} and \path{-20261004b} directories and
are superseded.

\section{Addendum: fresh potential-field re-solve restored}
At the user's direction, the fresh potential-field re-solve was restored as part
of the original algorithm, not as a fallback. A shifted problem that has no
primary point, has more than the unavoidable primary slack, or has no CLF
result is solved once more from a fresh rollout; ties retain the shift. The
other removed mechanisms stay removed: input-box enlargement, inherited budgets
and restoration, and replacement of an unusable shift. With the same source
otherwise, 656 of 656 tests pass (\path{tests-fresh-restart.csv}).

\begin{center}
\begin{tabular}{lccc}
\toprule
 & recovered & no solution & collision\\
\midrule
Exact observations, without re-solve & 11 & 3 & 0\\
Exact observations, with re-solve & 13 & 1 & 0\\
Noisy estimator, without re-solve & 1 & 13 & 0\\
Noisy estimator, with re-solve & 3 & 11 & 0\\
\bottomrule
\end{tabular}
\end{center}

With exact observations, every case recovers except 15-m/s brakingLead, which
has no primary point at hold 110. With the noisy estimator, 8-m/s headOn and
crossing and 15-m/s crossing recover. The other noisy stops follow the
constraint conflicts diagnosed above: brakingLead at frame~0, and the
short-horizon road and target rows. Per-case values are in
\path{campaign-metrics-fresh-restart.csv}; raw traces are in
\path{~/.cache/collisionAvoidance/fresh-restart-20261004/}.

\section{Addendum 2: re-solve only on trust-region infeasibility}
The user clarified the restored rule: the fresh rollout is used only when the
problem is infeasible inside the linearization trust region. The rule now
fires only when Clarabel certifies the shifted primary problem primal
infeasible. Positive slack is issued; numerical stalls and a CLF stage without
a result are reported. All 656 tests pass (\path{tests-infeasible-restart.csv}).
With exact observations 12 of 14 recover. 8 and 15~m/s brakingLead stop at
holds 26 and 110 with \texttt{clfNoNumericalResult}; the broader rule had
recovered the first and stopped the second. With the noisy estimator 3 of 14
recover, the same cases as before. No collision occurred
(\path{campaign-metrics-infeasible-restart.csv}).

\section{Addendum 3: trust-region expansion until feasible}
At the user's direction, a primary problem certified infeasible inside the
trust region is now re-solved on the same linearization with the trust scale
doubled until it becomes feasible, up to 16. If the shifted plan is still
infeasible at 16, the fresh rollout is expanded in the same way. The online
trust estimate is unchanged. All 657 tests pass.

With exact observations, 12 of 14 recover, as before. With the noisy estimator,
2 of 14 recover (previously 3), and there are no collisions. 15-m/s
turningCrossing now recovers; 8-m/s headOn and 15-m/s crossing no longer do.
Two new failure types appear. 8-m/s headOn hits an invalid tire operating point
at 3.95~s. 15-m/s curvedHeadOn has no usable anchor at 3.70~s: the previous
plan cannot be rolled out on the nonlinear model.

Before these failures the expansion reached 4 to 16 times the estimated scale.
The plan's model innovation grew from 0.08 to 1.7~m. Issued inputs jumped: the
braking ratio reached $-1.000$, and steering changed from 0.20 to $-0.82$~rad
in one hold. The expanded problems are feasible on a linearization that is no
longer accurate at that step (\path{campaign-metrics-trust-expansion.csv}).

\section{Addendum 4: expansion capped at the trust maximum}
The cap of 16 contradicted the linearization study: at a trust scale of 1 the
16-hold error is already 28~cm (p90, 15~m/s), and the tires saturate near 3.
The expansion is now capped at \texttt{trustMaximumScale} = 1 and the separate
parameter is removed. All 657 tests pass.

With exact observations, 12 of 14 recover; the result is identical without
expansion, with cap 16 and with cap 1. With the noisy estimator, 1 of 14
recovers, against 3 without expansion and 2 with cap 16; there are no
collisions. Six noisy runs stop between holds 73 and 79 (about 3.7~s). Each of
these stops becomes solvable only when the terminal road rows are removed.
When the target leaves $R$ the horizon shrinks to 8--29 nodes, and the ego is
still 1.8--5.7~m off the path. No trust expansion can satisfy the road rows
in that time (\path{noisy-drops-cap1.txt}). Expansion therefore does not
address the dominant noisy failure, and in the noisy runs it changed
trajectories that had recovered without it.

\section{Addendum 5: one road constraint at every node, four-lane road}
At the user's direction there is one road constraint, and it applies to the
whole prediction and the terminal alike: every ego rectangle corner stays
within \texttt{lateralClearance} at every predicted node after the measured
one, at every hold midpoint and at the endpoint. The 10-m centre corridor, its
configuration field, its test and its document are removed. Seed guidance is
clipped to the road less half the body width.

With the former $\pm4$~m scenario road this left too little room. In the
head-on fixture the ego centre must be at least 2.1~m from the target lane
centre, yet at most 3.05~m from the path for an aligned body, and about 2.5~m
at 0.24~rad yaw. The first head-on frame was infeasible; dropping either the
road rows or the collision rows solved it. Exact runs recovered 3 of 14 and
noisy runs 0 of 14, without collisions.

The scenario road now follows a four-lane road with 3.75~m lanes. The ego
follows the centre of the second lane from the right, so
\texttt{lateralClearance} = [right; left] = [5.625; 9.375]~m. With
this road, 638 of 638 tests pass. Exact runs recover 9 of 14; all five stops
are CLF-stage numerical stalls (\texttt{clfNoNumericalResult}). Noisy runs
recover 2 of 14. No collision occurs. In the \texttt{ode45} replay, two noisy runs on the
curved road touch or leave the road before they stop: 15~m/s curvedCrossing by
0.243~m and 15~m/s curvedHeadOn by 0.001~m. The road rows act on the estimated
state with corner coordinates linearized on the curved path; they do not bound
the realized plant
(\path{campaign-metrics-road-everywhere-four-lane.csv}).

\section{Addendum 6: US lane width and 10-ft shoulders}
The experiments are in the US, so the scenario road uses 12-ft (3.6576~m)
lanes. The ego still follows the centre of the second lane from the right. A
first run without shoulders (\texttt{lateralClearance} = [5.4864; 9.144]~m)
recovered 12 of 14 exact and 2 of 14 noisy runs, without collisions; 639 of
639 tests passed (\path{campaign-metrics-us-four-lane.csv}). Two noisy runs on
the curved road left the road, 15~m/s curvedHeadOn by 0.072~m and 15~m/s
curvedCrossing by 0.17~m.

With 10-ft (3.048~m) paved shoulders on both sides counted as road,
\texttt{lateralClearance} = [8.5344; 12.192]~m. All 639 tests pass. Exact runs
recover 11 of 14: 8 and 15~m/s brakingLead and 15~m/s turningCrossing stop
with \texttt{clfNoNumericalResult}. Noisy runs recover 4 of 14. Seven stop with
\texttt{pcbfNoNumericalResult}, 15~m/s brakingLead stops at frame~0 on the
solver time limit, and 15~m/s curvedCrossing ends with a tire-domain error. There are no collisions.
Three noisy 15-m/s runs (crossing, curvedHeadOn, curvedCrossing) leave even the
widened road, by 0.14--0.21~m, in the \texttt{ode45} replay
(\path{campaign-metrics-us-shoulder.csv}).

\section{Addendum 7: a road shoulder as a terminal set}
At the user's direction, driving onto a shoulder is also a terminal set; targets
are assumed never to enter a shoulder. The terminal set is now a union of three
modes: the target beyond~$R$ and separating, the ego rectangle on the right
shoulder, or on the left shoulder. A shoulder mode requires every corner inside
that shoulder ($\texttt{road.shoulderWidth}=[3.048; 3.048]$~m, the outer part of
\texttt{lateralClearance}) and the heading within 0.05~rad of the path; the
speed is free. Every mode keeps the road rows.

Every mode is solved and the issued mode has the least primary slack; modes
tied within \texttt{lexicographicTieTolerance} compete on CLF slack. A shifted
plan is one anchor for all modes; its extension follows path guidance toward
the region of the mode it ended in. At startup, and after every mode of a shift
is infeasible, each mode has its own fresh rollout: the potential field for
separation, and path guidance to the shoulder centre for a shoulder. On a
single shared anchor the shoulder modes were infeasible in all three startup
probes, because that anchor never approaches a shoulder. The trust scale is
expanded only while every mode is infeasible; per-mode expansion had cost eight
extra solves per frame. All 642 tests pass (\path{tests-shoulder-terminal.csv}).

\begin{center}
\begin{tabular}{lccc}
\toprule
 & recovered & no solution & collision\\
\midrule
Exact, shoulders as road only & 11 & 3 & 0\\
Exact, shoulder terminal & 11 & 3 & 0\\
Noisy, shoulders as road only & 4 & 10 & 0\\
Noisy, shoulder terminal & 3 & 11 & 0\\
\bottomrule
\end{tabular}
\end{center}

With exact observations the same three cases stop with
\texttt{clfNoNumericalResult}: 8 and 15~m/s brakingLead (now at holds 29 and
31) and 15~m/s turningCrossing at hold 5. A shoulder mode was issued in 9 of
14 exact runs, in six of them from the first frame, where all modes tie at zero slack and the
one-step CLF slack decides. With the noisy estimator, 8~m/s headOn,
acceleratingHeadOn and crossing recover. 8~m/s curvedHeadOn, previously
recovered, now stops at hold 1545 on a CLF numerical failure. 15~m/s brakingLead now runs 90 holds
instead of stopping at frame~0. No \texttt{ode45} replay leaves the road
(minimum margin 0.815~m); before, three noisy runs did. No run collides
(\path{campaign-metrics-shoulder-terminal.csv}, with per-mode hold counts).

The shoulder premise is not checked: a target predicted onto a shoulder is not
excluded by the shoulder mode, although the collision rows inside~$R$ remain.
Selection by one-step CLF slack among tied modes is a tie rule, not a measure
of how good a whole plan is.

\section{Addendum 8: shoulder terminal withdrawn, separating terminal set}
The user reported that a shoulder terminal had conflicted with the CLF in their
earlier work and never returned to nominal behaviour. The runs of Addendum~7
showed the same signs: recovery took longer in three exact runs (356 to 481,
159 to 196 and 154 to 218 holds), and two noisy curved head-on runs chose a
shoulder within eight holds of a CLF-stage stop. Constraint-drop replays of the
Addendum~6 failures showed that the perception-radius terminal was the sole
cause only in the noisy brakingLead cases at frame~0: following a lead never
leaves~$R$. The shoulder-terminal code is reverted.

At the user's direction the terminal set is relaxed to separation plus the
road. At the last node $(p-q)^\top(R(\psi)v-w)\ge0$; the distance condition is
removed. With a constant relative velocity the squared distance then has a
nonnegative first and second derivative, so the distance never decreases after
the endpoint. Because the endpoint can now be close to the target, it also
carries the collision rows. The row is linearized in position, yaw and body
velocity and divided by the anchor distance. The horizon ends at the first
anchor node whose separating speed is at least 0.5~m/s
(\texttt{terminalSeparatingMarginMetersPerSecond}); the row requires zero.

\begin{center}
\begin{tabular}{lccc}
\toprule
 & recovered & no solution & collision\\
\midrule
Exact, beyond $R$ and separating (Addendum 6) & 11 & 3 & 0\\
Exact, separating only & 14 & 0 & 0\\
Noisy, beyond $R$ and separating (Addendum 6) & 4 & 10 & 0\\
Noisy, separating only & 4 & 10 & 0\\
\bottomrule
\end{tabular}
\end{center}

With exact observations every case recovers, including both brakingLead runs,
and the subsequent median frame time is 3.5--11~ms. The longest horizon per exact run is
22--111 nodes. The noisy runs recovering are 8~m/s headOn,
acceleratingHeadOn, crossing and curvedHeadOn. The noisy brakingLead runs now
run 136 and 63 holds instead of stopping at frame~0. No \texttt{ode45} replay
collides or leaves the road; the minimum sampled gap is below the 0.20~m
planning margin in several runs (minimum 0.110~m).

Each noisy stop is a primal-infeasible primary problem. In the constraint-drop
replays the separating row is never the cause: removing it alone solves none of
the ten stops. Removing the collision rows solves seven, of which three also
solve with only the hard nominal collision rows removed; removing the road rows
solves six, three of them only that way
(\path{separating-terminal-drops.txt}, \path{campaign-metrics-separating-terminal.csv}).

The distance argument holds only for a constant relative velocity. A braking
lead, a turning target or the ego's own return to the path can close the
distance after the endpoint; no invariance is claimed.
\end{document}
